\let\negmedspace\undefined
\let\negthickspace\undefined
\documentclass[journal]{IEEEtran}
\usepackage[a5paper, margin=10mm, onecolumn]{geometry}
%\usepackage{lmodern} % Ensure lmodern is loaded for pdflatex
\usepackage{tfrupee} % Include tfrupee package

\setlength{\headheight}{1cm} % Set the height of the header box
\setlength{\headsep}{0mm}     % Set the distance between the header box and the top of the text

\usepackage{gvv-book}
\usepackage{gvv}
\usepackage{cite}
\usepackage{physics}
\usepackage{amsmath,amssymb,amsfonts,amsthm}
\usepackage{algorithmic}
\usepackage{graphicx}
\usepackage{textcomp}
\usepackage{xcolor}
\usepackage{txfonts}
\usepackage{listings}
\usepackage{enumitem}
\usepackage{mathtools}
\usepackage{gensymb}
\usepackage{comment}
\usepackage[breaklinks=true]{hyperref}
\usepackage{tkz-euclide} 
\usepackage{listings}
% \usepackage{gvv}                                        
\def\inputGnumericTable{}                                 
\usepackage[latin1]{inputenc}                                
\usepackage{color}                                            
\usepackage{array}                                            
\usepackage{longtable}                                       
\usepackage{calc}                                             
\usepackage{multirow}                                         
\usepackage{hhline}                                           
\usepackage{ifthen}                                           
\usepackage{lscape}
\begin{document}

\bibliographystyle{IEEEtran}

\title{12.36}
\author{EE25BTECH11022 - sankeerthan}
% \maketitle
% \newpage
% \bigskip
{\let\newpage\relax\maketitle}

\renewcommand{\thefigure}{\theenumi}
\renewcommand{\thetable}{\theenumi}
\setlength{\intextsep}{10pt} % Space between text and floats


\numberwithin{align}{enumi}
\numberwithin{figure}{enumi}
\renewcommand{\thetable}{\theenumi}

\textbf{Question:}Given the regression model:$Y_i = \alpha + \beta x_i + \epsilon_i$ with data:  \\ \begin{align}
\begin{tabular}{|c|c|}\hline
\text{Year} & \text{Production} \\ \hline
2001 & 30 \\ \hline
2002 & 35 \\ \hline
2003 & 36 \\ \hline
2004 & 32 \\ \hline
2005 & 37 \\ \hline
2006 & 40 \\ \hline
\end{tabular}
\end{align}

\textbf{Solution:}
Let
\begin{align}
x_i = \myvec{1 \\ 2 \\ 3 \\ 4 \\ 5 \\ 6}, Y = \myvec{30 \\ 35 \\ 36 \\ 32 \\ 37 \\ 40}
\end{align}

Let
\begin{align}
X = \myvec{1 & 1 \\1 & 2 \\1 & 3 \\1 & 4 \\1 & 5 \\1 & 6}
\end{align}

\begin{align}
X^T X =\myvec{6 & 21 \\21 & 91},X^T Y =\myvec{210 \\761} \\
\brak{X^T X}^{-1} =\myvec{0.87 & -0.20 \\-0.20 & 0.06} \\
\myvec{\alpha \\\beta}= \brak{X^T X}^{-1}\brak{X^T Y} \\
=\myvec{0.87 & -0.20 \\-0.20 & 0.06} \\
=\myvec{210 \\761}=\myvec{29.80 \\1.49}
\end{align}
Thus, $\alpha + \beta = 31.29$


\begin{figure}[h!]
    \centering
    \includegraphics[width=1\columnwidth]{figs/points.png}
    \caption{}
    \label{}
\end{figure}


\end{document}
